\documentclass{beamer}
\usepackage{nagoyatheme}

\usepackage[T1]{fontenc}
\usepackage[table,dvipsnames]{xcolor}
\usepackage{tabularx}
\usepackage{multirow}
\usepackage{amsmath, amssymb, amsthm}


\usepackage{tikz}
\usetikzlibrary{calc, arrows.meta, intersections, patterns, positioning, shapes.misc, fadings, through,decorations.pathreplacing}


\definecolor{ColorOne}{named}{nagoyaGreen}
\definecolor{ColorTwo}{named}{Gray}
\definecolor{ColorThree}{named}{Plum}


\title{G30 1st year report}
\author[WULFMAN]{Pierre-Emmanuel WULFMAN}
\institute[GSM, Nagoya U]{Graduate School of Mathematics}

\date{October 2024 - September 2025}


% For code snippets
\usepackage{listings}
\lstset{basicstyle=\ttfamily,columns=flexible,breaklines=true}

% Bibliography
\usepackage{csquotes}
\usepackage[style=authoryear]{biblatex}

\addbibresource{midmaster_presentation.bib}

\begin{document}

\begin{frame}[plain, noframenumbering]
   \titlepage
\end{frame}

\begin{frame}{Outline}
   \tableofcontents
\end{frame}

\section{Introduction}
\begin{frame}{Introduction}
   \begin{block}{}
      \begin{itemize}
         \item Brief self-introduction
         \item Goal of this presentation: highlight key learnings and experiences
      \end{itemize}
   \end{block}
   \begin{block}{}
      \tikzstyle{descript} = [text = black,align=center, minimum height=1.8cm, align=center, outer sep=0pt,font = \footnotesize]
      \tikzstyle{activity} =[align=center,outer sep=1pt]

      \begin{tikzpicture}[very thick, black]
         \small

         %% Coordinates
         \coordinate (O) at (0,0); % Origin
         \coordinate (P1) at (2.5,0);
         \coordinate (P2) at (5,0);
         \coordinate (P3) at (7.5,0);
         \coordinate (F) at (10,0); %End

         %% Filled regions
         \fill[color=ColorOne!20] rectangle (O) -- (P1) -- ($(P1)+(0,1)$) -- ($(O)+(0,1)$); % Studies
         % \path [pattern color=ColorOne, pattern=north east lines, line width = 1pt, very thick] rectangle ($(O)+(0.5,0)$) -- ($(O)+(10,0)$) -- ($(O)+(10,1)$) -- ($(O)+(0.5,1)$); % Something else
         \fill[color=ColorTwo!20] rectangle (P1) -- (P2) -- ($(P2)+(0,1)$) -- ($(P1)+(0,1)$); % Work
         \fill[color=ColorOne!20] rectangle (P2) -- (P3) -- ($(P3)+(0,1)$) -- ($(P2)+(0,1)$); % Current activity
         \fill[color=ColorTwo!20] rectangle (P3) -- (F) -- ($(F)+(0,1)$) -- ($(P3)+(0,1)$); % Future

         %% Text inside filled regions
         \draw ($(O)+(1.2,0.5)$) node[activity,ColorOne] {1 semester};
         \draw ($(P1)+(1.2,0.5)$) node[activity,ColorTwo] {2 semester};
         \draw ($(P2)+(1.2,0.5)$)  node[activity, ColorOne] {3rd semester};
         \draw ($(P3)+(1.2,0.5)$)  node[activity, ColorTwo] {4th semester};

         %% Description
         % \node[descript,fill=ColorTwo!15,text=ColorTwo](D2) at ($(P2)+(-2,-2.5)$) {%
         %    \textbf{Where}\\
         %    Project\\
         %    description};

         % \node[descript,fill=ColorThree!15,text=ColorThree](D3) at ($(P3)+(-2,-2.5)$) {%
         %    \textbf{Where}\\
         %    Project\\
         %    description};

         %% Events
         % \draw[<-,thick,color=black] ($(E1)+(0,0.1)$) -- ($(E1)+(0,1.5)$) node [above=0pt,align=center,black] {Unexpected\\event};
         % \draw [decorate,decoration={brace,amplitude=6pt}]($(E2)+(-1,1.2)$) -- ($(E2)+(0.5,1.2)$) node [black,midway,above=6pt] {Something};

         %% Arrows
         % \path[->,color=ColorTwo] ($(P2)+(-1,-0.1)$) edge [out=-90, in=130]  ($(D2)+(0,1)$);
         % \path[->,color=ColorThree]($(P3)+(-1,-0.1)$)  edge [out=-70, in=90]  ($(D3)+(0,1)$);

         %% Arrow
         \draw[->] (O) -- (F);
         %% Ticks
         \foreach \x in {0,2.5,5,7.5}
         \draw(\x cm,3pt) -- (\x cm,-3pt);
         %% Labels
         \foreach \i \j in {0.5/{2024/10},2.5/{2025/04},5/{2025/10},7.5/{2026/04},9.5/Graduation}{
         \draw (\i,0) node[below=3pt] {\j} ;
         }

      \end{tikzpicture}
   \end{block}
\end{frame}


\section{Seminars}
\begin{frame}{Quantum Computing Seminar (Le Gall's group)}
   \begin{itemize}
      \item Focused on theoretical models of quantum computing
      \item Presented a survey on the different models (quantum Turing machines, quantum circuits, etc.)
      \item Presented the adiabatic model and an implementation of TSP based on the work of ~\parencite{Kieu_2019}\footfullcite{Kieu_2019}
   \end{itemize}
\end{frame}

\begin{frame}{Type Theory and Formal Proving (Garrigue's group)}
   \begin{itemize}
      \item 1st semester: \citetitle{Girard_1989}\footfullcite{Girard_1989} cover to cover
      \item 2nd semester:\\
            \citetitle{Reynolds_1998}\footfullcite{Reynolds_1998} Chap. 1-3, present. 2.4-2.8, 3.7-3.8.
            \\
            \citetitle{Appel_Dockins_Hobor_Beringer_Dodds_Stewart_Blazy_Leroy_2014}\footfullcite{Appel_Dockins_Hobor_Beringer_Dodds_Stewart_Blazy_Leroy_2014}
            Chap 1-21, present. 4-5,9-10,14-15,20-21.
   \end{itemize}
\end{frame}

\section{Separation Logic}

%------------------------------------------------
\begin{frame}{Motivation}
   \begin{itemize}
      \item  \textbf{Program verification}: Ensure programs respect specification.
      \item \textbf{specification}:
            \begin{block}{Hoare Triple}
               \centering
               \begin{tikzpicture}[every node/.style={align=center}]
                  \tikzstyle{descript} = [text = black,align=center, align=center, outer sep=0pt,font = \footnotesize]
                  % Nodes
                  \node (pre) at (4,0) {$\{P\}$};
                  \node (cmd) at (5,0) {$C$};
                  \node (post) at (6,0) {$\{Q\}$};
                  %% Description
                  \node[descript,fill=ColorTwo!15,text=ColorOne](D1) at ($(pre) + (-1,+1)$) {{Precondition}};
                  \node[descript,fill=ColorTwo!15,text=ColorOne](D2) at ($(post) + (1,1)$) {{Postcondition}};
                  \node[descript,fill=ColorTwo!15,text=ColorOne](D3) at ($(cmd) + (0,-1)$) {{Program}};

                  %% Arrows
                  \path[->,color=ColorTwo] ($(pre) + (0,0.3)$) edge [out=90, in=-90]  ($(D1) + (0,-0.2)$);
                  \path[->,color=ColorTwo] ($(post) + (0,0.3)$) edge [out=90, in=-90]  ($(D2) + (0,-0.2)$);
                  \path[->,color=ColorTwo] ($(cmd) + (0,-0.2)$) edge [out=-130, in=60]  ($(D3)+(0,0.2)$);
               \end{tikzpicture}
            \end{block}
      \item \textbf{Hoare Logic}: Axiom system and inference rules.
   \end{itemize}
\end{frame}

%------------------------------------------------
\begin{frame}{Example and Limitations}
   \begin{itemize}
      \begin{block}{Variable Assignment}
         \[
            \{a = 0\} \ a := 42 \ \{a = 42\}
         \]
      \end{block}
      \item trivial

            \pause
            \vspace{0.5em}
            \begin{block}{Array Assignment}
               \[
                  \{a[i] = 0 \land a[j] = 0\} \ a[i] := 42 \ \{a[i] = 42 \land a[j] = 0\}
               \]
            \end{block}
      \item undecidable with aliasing (i.e., $i = j$?)
   \end{itemize}
\end{frame}

%------------------------------------------------
\begin{frame}{Solution - Separation Logic}
   \begin{itemize}
      \item Extends Hoare logic with the \textbf{separating conjunction}:
            \[
               P * Q
            \]
      \item Meaning: $P$ and $Q$ are predicates on \textbf{disjoint} parts of memory.
      \item Enables \textbf{local reasoning}: reason only about relevant memory.
   \end{itemize}

   % \pause
   \begin{center}
      \begin{tikzpicture}[scale=0.8]
         \tikzstyle{descript} = [text = black,align=center, align=center, outer sep=0pt,font = \footnotesize]
         % Nodes
         \node (pre) at (1,0) {$\{a[i] = 0 * a[j] = 0\}$};
         \node (cmd) at (5,0) {$a[i] := 42$};
         \node (post) at (9,0) {$\{a[i] = 42 * a[j] = 0\}$};
         %% Description
         % \node[descript,fill=ColorTwo!15,text=ColorOne](D1) at ($(pre) + (-1,+1)$) {{Precondition}};
         \node[descript,fill=ColorTwo!15,text=ColorOne](D2) at ($(post) + (-2,-1.2)$) {Proof of right heaplet is trivial};
         \node[descript,fill=ColorTwo!15,text=ColorOne](D3) at ($(pre) + (-1,-1.2)$) {Only touch left heaplet};

         %% Arrows
         % \path[->,color=ColorTwo] ($(pre) + (0,0.3)$) edge [out=90, in=-90]  ($(D1) + (0,-0.2)$);
         \path[->,color=ColorTwo] ($(post) + (1.2,-0.2)$) edge [out=-130, in=60] ($(D2) + (0,0.3)$);
         \path[->,color=ColorTwo] ($(cmd) + (0,-0.2)$) edge [out=-130, in=60] ($(D3)+(0,0.3)$);
      \end{tikzpicture}
   \end{center}
\end{frame}

\begin{frame} {Generalization - Separation Algebra}
   \begin{block}{Definition}
      \begin{itemize}
         \item $A$ is a type $\simeq$ set (e.g., memory cells),
         \item $\oplus$\footnote{we write $\oplus(a,b,c)$ as $a \oplus b = c$} partial tertiary operation on $\mathcal{P}(A)$ (the powerset of $A$)
         \item commutative, associative, positive $\to$ permission algebra
         \item with unit $\to$ separation algebra
      \end{itemize}
   \end{block}
   \begin{block}{Examples}
      \begin{table}
         \begin{tabular}{c c c }
            name     & type $A$     & join operation                         \\
            % Heaplets & $\mathbb{N}\to \mathbb{N}$   &                       \\\hline
            Booleans & bool         & $(c = a \lor b) \land \neg(a \land b)$ \\
            Counters & $\mathbb{N}$ & $c =\max(a,b)$                         \\
            Discrete & any set      & $a = b = c$                            \\
            % Rational shares & $\mathbb{Q} $ & $a + c = b$                            \\
         \end{tabular}
      \end{table}
   \end{block}
\end{frame}


\begin{frame}{Make your own logic!}
   \begin{block}{operation}
      \begin{itemize}
         \item Prod $<A, J_A>$ and $<B, J_B>$ $\to$ $<A\times B, J>$
         \item Func $<B, J_B>$ and $A$ $\to$ $<A\to B, J>$
         \item Indexed Prod/Sum, etc.
      \end{itemize}
   \end{block}
   \begin{block}{Simple imperative program}
      \begin{itemize}
         \item State : <stack : var $\to$ option adr, heap : adr $\to$ option adr>
         \item Predicate : share \emph{stack}, disjoint \emph{heap}.
         \item Algebra :  Equiv(var $\to$ option adr )$\times $ adr $\to $ Lower ( Discret (option adr))
      \end{itemize}
   \end{block}
\end{frame}
%------------------------------------------------
% \begin{frame}{Example}
%    %    \begin{lstlisting}[language=C]
%    % int *p = malloc(sizeof(int));
%    %%%p = 42;
%    % free(p);
%    %   \end{lstlisting}

%    \vspace{0.5em}
%    \begin{block}{Hoare Triple}
%       \[
%          \{ emp \} \quad p = malloc(...); *p = 42; free(p); \quad \{ emp \}
%       \]
%    \end{block}

%    \pause
%    \begin{itemize}
%       \item \texttt{emp} means empty heap.
%       \item Allocation, update, and deallocation are localized.
%       \item No interference with unrelated memory.
%    \end{itemize}
% \end{frame}

%------------------------------------------------
\begin{frame}{Applications}
   \begin{itemize}
      \item Static analysis tools (Rust Type System)
      \item Modeling concurrent algorithms (\emph{Concurrent Separation Logic})
      \item Proof assistants
   \end{itemize}
   \vspace{0.5em}
   \pause
   \begin{block}{Outlook}
      Research continues with automated tools (e.g. \emph{Infer}, \emph{Iris})
   \end{block}
\end{frame}


% \section{Coursework}
% \begin{frame}{Representation Theory}
%    \begin{itemize}
%       \item Core mathematics course
%       \item Focus on algebraic structures and their applications
%       \item Reinforced mathematical maturity and abstract thinking
%    \end{itemize}
% \end{frame}

% \begin{frame}{Quantum Information Theory}
%    \begin{itemize}
%       \item Learned quantum states, channels, entropy
%       \item Applications in communication and computation
%       \item Connection with research interests
%    \end{itemize}
% \end{frame}

\section{Other Activities}
\begin{frame}{}
   \begin{block}{Teaching}
      \begin{itemize}
         \item QTA for [G30] Data Science
      \end{itemize}
   \end{block}

   \begin{block}{Conferences}
      \begin{itemize}
         \item Speaker at the Rocq Workshop 2025 (Rejkavik)
      \end{itemize}
      Attended :
      \begin{itemize}
         \item TPP 2024 (Kyushu)
         \item PPL 2025 (Gamagori)
         \item QI2025 Satellite Workshop (Osaka)
      \end{itemize}
   \end{block}
\end{frame}


\section{Conclusion}
\begin{frame}{Master's Program Progress}
   % \begin{itemize}
   %    \item Completed 1st year of Master's program
   %    \item Courses in quantum computing and mathematics
   %    \item Preparing for research project in 2nd year
   % \end{itemize}
   \begin{block}{credits}
      \begin{table}[!h]
         \begin{tabularx}{\textwidth}{|X|c|c|}
            \hline
            \rowcolor[HTML]{006E4F}
            {\color[HTML]{FFFFFF} type}  & {\color[HTML]{FFFFFF} obtained} & {\color[HTML]{FFFFFF} required} \\\hline
            MC Basic Specialized Courses & 2                               & 2                               \\\hline
            MC Specialized Courses       & 8                               & 8                               \\\hline
            Seminar                      & 10                              & 20                              \\\hline
            Optional Courses ?           & 5                               & ?                               \\\hline
         \end{tabularx}
      \end{table}
   \end{block}
   \begin{block}{Master Thesis}
      \begin{itemize}
         \item Status : Looking for a problematic
         \item Supervisor: Prof. Garrigue
      \end{itemize}
   \end{block}
\end{frame}

\begin{frame}
   \centering \Huge Questions?
\end{frame}

\appendix


\begin{frame}{References}
   \printbibliography
\end{frame}

\end{document}
